\chapter{Photosynthesis: Sunlight Is Not Directly Made into Food}

\begin{kaichang}
Before swallowing this morning's bread, rice, or steamed bun, think: it was once a wheat or rice plant. Earlier still, that plant grew from a seed weighing a few grams, yet a mature wheat plant weighs hundreds of times more.

Where did all that additional mass come from?

Guess and write your answer down. There are four candidates: eaten from the soil, taken from water, made by turning sunlight into a body, or drawn from air. I suspect many will choose soil; plants have roots in it, after all. Others may choose sunlight: do textbooks not say plants use sunlight to make food?

We will check your answer at the end. The answer is among those choices, but works quite differently from what most people imagine. Sunlight is indeed a plant's lifeline, but it is wages, not raw material. A plant does not, and cannot, turn light directly into its body.
\end{kaichang}

\section{A Willow Tree and the Mystery of Mass}

Start with a real experiment. In the 1640s, the Belgian physician Jan Baptista van Helmont filled a large wooden tub with dried soil weighing \ud{90}{kg}, planted a \ud{2.3}{kg} willow sapling, and covered the soil to keep dust out. For five years he added only rainwater and distilled water. The willow grew to more than \ud{76}{kg}, while the soil, dried and weighed again, had lost only about \ud{60}{g}.

Over seventy kilograms of new wood and foliage, with just sixty grams supplied by soil. Van Helmont concluded that a plant's body came almost entirely from water. He was half wrong: water is one raw material, but he missed the largest source of mass, air. His era lacked the idea that gases have mass and participate in chemical reactions, so we cannot blame him. Yet the experiment's quantitative conclusion still holds: \textbf{soil is not the main source of a plant's body}.

What about sunlight? Light is energy, not matter. You cannot weigh it and stitch it into wood; energy has no entry in the mass ledger, at least on everyday chemical scales, where the mass change involved in photosynthesis is too small to measure. If a plant is neither made from soil nor transformed light, the remaining sources are water and air. Later we will see that more than ninety percent of a leaf's dry mass ultimately comes from carbon dioxide and water. Your exhaled waste gas supplies a plant's building bricks.

Another observation awaits by a pond: at noon on a sunny day, little bubbles coat aquatic leaves and rise in streams. What is in them? Where does it come from? Why are there almost none on cloudy days? Keep these questions as this chapter's signposts.

\section{Zooming In: Green Factories in Leaves}

Zoom inside a leaf. Its cells contain tiny green bodies called \textbf{chloroplasts}, just a few micrometers across, with dozens in one mesophyll cell. Chapter 51 described cell membranes as boundaries formed by phospholipid bilayers. Chloroplasts go further: they have two envelope membranes and internal stacks of flattened coins called thylakoids. These coin membranes are densely embedded with two kinds of components: light-absorbing pigments and protein machines.

Where does green come from? The leading thylakoid pigment is chlorophyll. As Chapter 10 noted, its molecule holds a magnesium atom at its center, surrounded by a carbon--nitrogen ring and trailing a long hydrophobic tail that anchors it in the membrane. Chlorophyll is choosy about light: it strongly absorbs blue-violet light around $430$ nanometers and red light around $660$ nanometers, while largely ignoring green light around $500$--$600$ nanometers. Most green light is reflected or transmitted, making leaves look green. \textbf{A leaf's color reveals the light it does not consume}.

Why particular wavelengths? Return to Chapter 9: light comes in packets, or photons, whose energy depends on wavelength. Shorter wavelengths mean more energy per packet. Molecular electrons accept only energies matching their allowed steps. Blue-violet and red photons match particular electronic transitions in chlorophyll, so absorption kicks an electron to a higher energy level. Green light does not match and passes through. Like a vending machine accepting particular coins, this is structural matching, not preference.

An excited electron is unstable. In many materials, it quickly falls back, returning energy as heat or faint fluorescence; sunscreen molecules dispose of ultraviolet energy this way. Chloroplasts do something different: \textbf{they intercept the electron before it falls back} and send it into an electron-transport chain. Photosynthesis truly begins not with absorbing sunlight, but with using sunlight to lift electrons from low to high energy and capture them in molecular pockets.

\begin{shiyan}
\textbf{Observation: Aquatic Plants Bubble in Light}

\textbf{Materials:} a sprig of hornwort from an aquarium or plant shop, or a fresh pothos cutting, which produces fewer bubbles; a transparent glass, clean water, a desk lamp, scissors, and a timer.

\textbf{Translator safety note:} Pothos is listed as poisonous by \href{https://www.poison.org/articles/plant}{Poison Control}; it is not a food ingredient. Have an adult identify and cut the plant, keep it away from mouths and eyes, wash hands afterward, and keep plants and experimental water away from children and pets.

\textbf{Procedure:} (1) Cut an approximately 10-centimeter piece of aquatic plant underwater to keep air out of the cut. (2) Place it cut-end upward in the glass and cover with water. (3) Position the lamp about 20 centimeters away and count bubbles from the cut end for 5 minutes. (4) Switch the lamp off or move the glass into darkness and count for another 5 minutes.

\textbf{Observations:} bubbles increase markedly in light and almost stop in darkness. They emerge from the single cut, not randomly across all the leaves. This shows that gas is produced inside the plant, rather than dissolved air escaping as water warms.

\textbf{Safety:} the experiment itself uses no dangerous chemicals, but a lamp shade becomes hot after prolonged use; do not touch it. Keep the glass away from the lamp and avoid splashing electrical equipment. Identifying the gas requires chemical testing: do that only with a teacher in a school laboratory, never by collecting and igniting gas yourself.
\end{shiyan}

\section{A Relay: From Water to ATP}

Zoom closer to see what the captured high-energy electrons do. The process is a relay on and near the thylakoid membrane, divided into two stages. The first directly requires light and is called the \textbf{light reactions}; the second does not directly require light and is called the \textbf{carbon reactions}. Textbooks often call these dark reactions, but that misleading name suggests a restriction that does not exist: they run day and night, so we will not use it.

First leg: split water. Losing an electron leaves chlorophyll electron-deficient; it must immediately replace the electron or the machine stops. Water supplies the replacement. A manganese-containing protein machine in the thylakoid membrane forcibly extracts electrons from water and gives them to chlorophyll. Water's hydrogen remains inside as hydrogen ions, \chem{H^+}; the oxygen atoms from two water molecules pair up into \chem{O_2} and are released. This is the source of the bubbles at your plant cutting.

Pause for three seconds over this chapter's first firm fact: \textbf{the oxygen plants release comes from water, not carbon dioxide}. Water is split; carbon dioxide does not contribute a single oxygen atom to the air throughout the process. This is not a guess but a conclusion firmly supported by isotope experiments, whose story comes later.

Second leg: electrons descend while protons are pumped upward. High-energy electrons from water are not immediately stored. They pass stepwise through a sequence of membrane proteins, losing a little energy at each transfer, almost the same script as the mitochondrial electron-transport chain in Chapter 53. That energy pumps \chem{H^+} from outside the membrane into the thylakoid lumen, building a concentrated, acidic proton gradient. Protons then flow back downhill through a nanoscale rotary motor, ATP synthase, twisting ADP into ATP. Chemiosmosis, introduced in Chapter 53, appears again. Evolution uses the same trick in two systems: one extracts energy by dismantling sugar, the other charges up using light.

Third leg: store electrons in a readily accessible account. At the chain's end, electrons and \chem{H^+} pass to the carrier \chem{NADP^+}, producing NADPH. NADPH is a close relative of Chapter 50's NADH, another cash van carrying high-energy electrons.

The light reactions are complete. Take inventory: oxygen, a released by-product; ATP, energy cash; and NADPH, a high-energy electron account. Notice that \textbf{not one carbon atom has yet been fixed}. Light has finished its job. What it bought was not sugar but two high-energy currencies.

\section{Carbon Reactions: Buying Carbon with Energy Currency}

The next stage occurs in the chloroplast's fluid stroma without directly requiring light. Plants spend newly earned ATP and NADPH to weld atmospheric carbon dioxide into organic compounds. This circular assembly line, the \textbf{Calvin cycle}, has roughly three stages.

First, capture carbon. Atmospheric \chem{CO_2} is dilute, only about four parts in ten thousand. A chloroplast enzyme captures a \chem{CO_2}, attaches it to a five-carbon sugar, and immediately splits the product into two three-carbon molecules. The enzyme is Rubisco, an awkward name worth remembering because it is Earth's most abundant protein: the world's Rubisco totals hundreds of millions of tonnes. Why such lavish investment? It works slowly, capturing only a few \chem{CO_2} per second, versus hundreds or thousands of substrates for typical enzymes. Plants compensate with quantity. A substantial share of every patch of green you see is actually Rubisco's color.

Second, reduction. The three-carbon molecules successively spend ATP's chemical energy and NADPH's high-energy electrons to become a three-carbon sugar. This is where the light reactions' cash and savings are spent, connecting the two stages: light reactions supply currency; carbon reactions consume it.

Third, keep the cycle running. After several turns capturing 3 \chem{CO_2}, the net gain is 1 three-carbon sugar that leaves the cycle. The remaining three-carbon molecules spend more ATP to regenerate five-carbon sugars for another round of capture. Exported three-carbon sugars combine into glucose, which is polymerized into stored starch as in Chapter 46 or used to build cellulose, proteins, and fats.

Return to the title: plants do not turn sunlight directly into food. The actual causal chain is that light energy splits water, releasing oxygen and lifting electrons to high energy; the electrons' energy is exchanged in portions for ATP and NADPH; those currencies drive enzymes that join and reduce \chem{CO_2} into sugar. \textbf{Sugar's substance comes from carbon dioxide and water; its energy is sunlight's legacy}. Light is the electricity for the factory, not the bricks for its buildings.

\section{Symbols: Why the Equation Needs Twelve Waters}

At last, we can write the reaction. Many textbooks abbreviate photosynthesis as:

\[\chem{6CO_2} + \chem{6H_2O} \rightarrow \chem{C_6H_{12}O_6} + \chem{6O_2}\]

This is not wrong, but it hides a key fact. Count oxygen atoms: the left has 12 in \chem{CO_2} and 6 in water, totaling 18; on the right, glucose takes 6 and \chem{O_2} takes 12. Could six waters apparently supply oxygen for six \chem{O_2}? That is the trap in the abbreviated version. Since all the oxygen gas comes from water, water must appear on the left \textbf{twice}: once as the electron donor being split and once as the reaction environment. A more revealing equation is:

\[\chem{6CO_2} + \chem{12H_2O} \rightarrow \chem{C_6H_{12}O_6} + \chem{6O_2} + \chem{6H_2O}\]

All 12 oxygen atoms in the left-hand waters enter the right-hand \chem{6O_2}; all the oxygen in glucose and the newly appearing 6 waters comes from \chem{CO_2}. The actual water-splitting step in the light reactions is:

\[\chem{2H_2O} \rightarrow \chem{O_2} + \chem{4H^+} + \chem{4e^-}\]

Remember this asymmetry: releasing 1 \chem{O_2} requires splitting 2 waters and removing 4 electrons. Those electrons are the whole assembly line's original capital.

\begin{zhengju}
How do we know oxygen comes from water rather than carbon dioxide? Both contain identical oxygen atoms, indistinguishable in an ordinary experiment. It is like mixing two identical glasses of water: how can you prove which glass a molecule came from?

Isotopes provide labels. Chapter 7 explained that an element can have versions with the same proton number but different neutron numbers. Ordinary oxygen is oxygen-16; stable oxygen-18 has two extra neutrons and nearly identical chemistry but can be distinguished instrumentally. In the 1930s, Samuel Ruben and Martin Kamen's team at Stanford University devised two culture mixtures: ordinary water plus oxygen-18-labeled \chem{CO_2}, and oxygen-18-labeled water plus ordinary \chem{CO_2}. They let green algae photosynthesize and tested the released \chem{O_2} for the label.

The result was clear: labeled \chem{CO_2} produced ordinary oxygen gas; labeled water produced oxygen-18-containing gas. Every oxygen atom's origin was registered to water. Published in 1941, these experiments became a cornerstone of photosynthesis research.

Someone had already guessed the answer by another route. The Dutch-born microbiologist C. B. van Niel studied purple sulfur bacteria, which photosynthesize but \textbf{do not release oxygen}. They split hydrogen sulfide, \chem{H_2S}, rather than water, producing sulfur granules instead of oxygen. Van Niel saw the parallel: bacteria split \chem{H_2S} and release sulfur; plants split \chem{H_2O} and release oxygen. The script is identical, with different starting molecules and residues. He inferred that plant oxygen must also come from the molecule being split: water. An analogy from an unusual organism was confirmed by isotope experiments a decade later. Many major scientific conclusions begin with conjecture and receive evidence later, but evidence always decides whether they stand.
\end{zhengju}

\section{Photosynthesis and Respiration: Different Routes}

Now address a misconception nourished by countless abbreviated equations. Put photosynthesis beside Chapter 53's cellular respiration:

\begin{itemize}[itemsep=2pt]
\item Photosynthesis: \chem{6CO_2} + \chem{12H_2O} $\rightarrow$ \chem{C_6H_{12}O_6} + \chem{6O_2} + \chem{6H_2O}
\item Cellular respiration: \chem{C_6H_{12}O_6} + \chem{6O_2} $\rightarrow$ \chem{6CO_2} + \chem{6H_2O}
\end{itemize}

They look remarkably alike. Reverse the arrow and they are almost mirror images. Hence the widespread claim that photosynthesis is respiration in reverse.

\begin{wuqu}
``Photosynthesis is respiration running backward; plants photosynthesize by day and respire at night.''

Both halves are wrong. First, \textbf{the routes are completely different}. Reversing a net equation is easy, but cellular chemistry proceeds stepwise, and one wrong step gives a different substance. Respiration uses glycolysis, the citric-acid cycle, and electron transport in the cytoplasm and mitochondria, as in Chapters 52 and 53. Photosynthesis uses photosystems and the Calvin cycle in chloroplasts, with almost no overlap between the assembly lines' enzymes. High-speed rail and air travel from Beijing to Shanghai share endpoints, but flying is not riding a train backward. Even the atom routes differ: oxygen in respiratory \chem{CO_2} comes from glucose and water; oxygen in photosynthetic \chem{O_2} comes from water. Isotope tracing identifies each route clearly: these are not the same atoms simply shuffled back and forth.

Second, \textbf{plants respire both day and night}. Leaves make sugar, but leaf cells must also eat. Maintaining membranes, making proteins, and repairing damage all cost ATP, mostly still supplied by mitochondria burning sugar. Measured daytime oxygen release merely means photosynthetic production exceeds respiratory consumption: the net balance of \textbf{two assembly lines working simultaneously}. At night, photosynthesis stops while respiration continues, and plants consume oxygen and release carbon dioxide like animals. Is the pothos in your bedroom stealing oxygen at night? Do not worry: one plant's respiratory demand is far smaller than that of another sleeping person, or even a cat.
\end{wuqu}

The two lines' deeper connection lies not in symmetric equations but in \textbf{shared molecular machinery}: both have electron-transport chains, store potential energy in proton gradients, and collect it through ATP synthase. Chapter 53's chemiosmosis and this chapter's light reactions are like the same engine in two vehicles. This is evolutionary reuse, not coincidence: the common ancestors of photosynthetic and respiratory bacteria long ago invented membrane proton gradients for storing and retrieving energy. Part X's evolution chapters will explore that ancestry; keep it in mind for now.

\section{One Recipe, Three Ways of Life}

The Calvin cycle is the carbon-assembly workshop of almost all green plants, but its operating environments vary widely. The central difficulty is Rubisco's occupational flaw: it cannot distinguish \chem{CO_2} from \chem{O_2}. Its carbon-capturing active site sometimes binds oxygen instead, producing an unusable intermediate that costs ATP to recover and reprocess. This is photorespiration. Higher temperatures and tightly closed stomata conserving water lower internal \chem{CO_2} and raise \chem{O_2}, increasing mistakes. At midday in hot, dry summers, plants such as rice and wheat can lose a substantial portion of fixed carbon this way, like a factory producing goods while reworking defects.

Evolution has produced three responses to this problem:

\begin{itemize}[itemsep=2pt]
\item \textbf{C$_3$ plants}, including rice, wheat, soybeans, and most trees, use the Calvin cycle directly. Their first stable fixation product has three carbons, hence C$_3$. They are most economical in cool, moist conditions, requiring no extra infrastructure, but photorespiration burdens them in hot, dry weather.
\item \textbf{C$_4$ plants}, including maize, sugarcane, and sorghum, add a preliminary pump. A less choosy enzyme first captures \chem{CO_2} in a four-carbon molecule in mesophyll cells, hence C$_4$. The molecule travels deeper into bundle-sheath cells to unload, raising local \chem{CO_2} severalfold. Surrounded by carbon dioxide, Rubisco rarely makes mistakes. Each fixed carbon costs extra ATP, the concentration pump's electricity bill, but under strong light and heat the expense pays off.
\item \textbf{CAM plants}, including cacti, pineapples, and agaves, separate stages in time rather than space. By day their stomata stay shut, letting out no water; on cool nights they open wide and store incoming \chem{CO_2} as malic acid. After dawn they close up and slowly release the stored carbon for the Calvin cycle. Water conservation is extreme, at the cost of slow growth. Desert plants are in no hurry.
\end{itemize}

None is superior in every setting; each matches its environment. This is natural selection's typical approach: not perfect design, but modifications to available parts until they work well enough. Rubisco's flaw remains unfixed, not for lack of desire, but because every intermediate step toward a fix is worse than the present state. Evolution cannot cross that valley.

\section{Where This Model Breaks Down}

This chapter's picture is useful, but remember three limits.

First, \textbf{photosynthesis does not necessarily release oxygen}. Van Niel's purple sulfur bacteria have photosynthetic pigments, electron transport, and carbon fixation, but split \chem{H_2S} rather than water and release no oxygen. Oxygenic, water-splitting photosynthesis is only one branch of a large family. It happens to be the branch that transformed Earth: today's atmosphere is one-fifth oxygen, all its industrial waste from billions of years of water splitting. The Great Oxidation Event returns in Chapter 56's account of early life.

Second, \textbf{net equations hide rates}. ``Enough light, water, and carbon dioxide will produce sugar'' balances the books but says nothing about speed. Production in a real leaf is limited by its slowest step: light reactions in dim light, sluggish Rubisco under strong light, and closing stomata at high temperatures. Greenhouse agriculture follows this logic, supplementing \chem{CO_2} from about four to eight parts per ten thousand and extending artificial illumination. Changing the bottleneck can significantly raise yields: the chapter's model becomes agricultural technology.

Third, \textbf{efficiency is much lower than you might expect}. Usually only one to a few percent of sunlight falling on leaves ends up stored in sugar. Green and infrared light are not absorbed; some absorbed energy escapes as heat; photorespiration consumes another portion. Before criticizing plants' inefficiency, remember that they have run on free light, air, and water for over two billion years at zero fuel cost, supporting the entire biosphere along the way. Human photovoltaic panels are more efficient, but their manufacturing and electricity accounts are hardly cheap.

\begin{qiaoliang}
How much sugar and oxygen can one square meter of vigorous foliage produce daily? A typical measured net photosynthetic rate fixes $15$ micromoles, $15\times 10^{-6}$ mol, of \chem{CO_2} per square meter per second. Assume $10$ effective hours of strong summer sunlight:

\[
15\times 10^{-6} \times 3600 \times 10 \approx 0.54\ \text{mol CO}_2/\text{day}
\]

The net equation converts every 6 \chem{CO_2} into 1 glucose and releases 6 \chem{O_2}. That means $0.09$ mol of glucose daily, about \ud{16}{g}, plus $0.54$ mol of oxygen. At roughly $24$ liters per mole at room temperature, that is about $13$ liters.

Compare yourself: a resting secondary-school student consumes a little over $500$ liters of oxygen daily. Thus \textbf{offsetting one day's breathing requires about 40 square meters of vigorous leaves working at full capacity in sunlight for a day}. A football-field-sized lawn balances roughly twenty people's breathing, before allowing for the grass's own nighttime respiration. Next time you hear ``the city's green lungs,'' remember the concrete atomic account behind it.
\end{qiaoliang}

\section{Back to the Bread}

Time to settle the opening question: where does the mass that makes a wheat plant hundreds of times heavier than its seed come from?

Not soil: van Helmont's tub rules that out, except for a small mineral contribution. Not direct conversion of sunlight: light becomes no atom. Chlorophyll captures it, exchanges it for ATP and NADPH, and the energy is spent and dissipated. The answer is \textbf{mostly atmospheric carbon dioxide, together with water}. Carbon is the brick, water the mortar, and light the electricity. Your bread is essentially solidified air plus captured sunlight. A writing teacher might mark down the phrase, but the atomic accounts support it exactly.

The bubbles from the aquatic-plant cutting are explained too. Oxygen released by water splitting travels through the plant's air spaces and escapes at the wound. The lamp pays the wages; switch it off, and water splitting stops and bubbles cease.

One question may have been on your mind: if plants turn \chem{CO_2} into sugar, why not copy them with an artificial device to solve food and climate problems? Scientists are trying: it is called artificial photosynthesis. The obstacle is not the principle, already covered here, but engineering. Water is among the hardest stable molecules to split. To do it cheaply, durably, and efficiently, today's human catalysts still fall far short of the robust manganese machine in chloroplasts. Nature's solution, refined for over two billion years, is not easy to copy.

\section{What Happens After Sugar Is Made?}

This chapter followed glucose only as far as the factory exit. A living cell cannot simply stockpile starch: sugar must be burned for ATP in Chapters 52 and 53, made into cellulose walls, converted into seed oils, and supply carbon skeletons for proteins and nucleic acids. The green factory's workshops are connected not by one-way streets but by a mutually regulating network. How insulin directs whole-body fuel allocation, and what the body burns during fasting, both belong to this network. Next chapter, we zoom out from individual assembly lines to metabolism as a whole.

\begin{tiaozhan}
How expensive is carbon fixation? For a net gain of 1 three-carbon sugar, the Calvin cycle fixes 3 \chem{CO_2} and spends 9 ATP and 6 NADPH. How many photons buy that currency? Every 4 electrons extracted from water, releasing 1 \chem{O_2}, must each be kicked upward by two successive photosystems, requiring at least $8$ photons. Their flow yields roughly 3 ATP and 2 NADPH. Scaling the account, the currency needed to fix 3 \chem{CO_2} corresponds to around $30$ photons. Those 30 red photons carry about $5100$ kilojoules per mole, while three carbons stored as sugar retain about $1400$ kilojoules, giving a combined on-paper efficiency near thirty percent. Wasteful? Remember that this is almost the ceiling for all biochemical processes: every energy handoff loses some as heat, the unavoidable toll of the second law of thermodynamics. Skipping this box does not affect the main argument.
\end{tiaozhan}

\section*{Practice}

\begin{sikaoti}
\item Draw a material-flow diagram using boxes and arrows to trace \chem{CO_2}, water, light energy, oxygen, ATP, NADPH, and glucose through photosynthesis. Mark each step's chloroplast location, thylakoid membrane or stroma. Note beside the diagram that boxes and arrows are human representations.
\item A student grows green algae in ordinary water with oxygen-18-labeled \chem{CO_2}. The label appears in glucose and newly formed water, but not released \chem{O_2}. Does this support or refute oxygen coming from water? Explain the reasoning.
\item Seal a green plant in a transparent glass enclosure with sufficient light. Over several days, \chem{CO_2} falls and \chem{O_2} rises, eventually reaching stable levels. Explain this using simultaneous photosynthesis and respiration and their net balance. Predict what happens to the gases if the enclosure is moved into complete darkness.
\item Maize, a C$_4$ plant, and rice, a C$_3$ plant, grow in the same hot, dry, strongly sunlit field. Which is more likely to yield more, and why? Would your prediction change in cool, moist conditions?
\item Use this chapter's statement ``light is electricity, not bricks'' to refute: ``More sunlight makes plants stronger, so plants grow by eating sunlight.'' Include at least three molecular or structural names from the chapter.
\item Estimate: a pothos on your balcony has about 0.1 square meters of leaf area. Using the bridge box's account, estimate its oxygen output during a vigorous day's growth. Explain why this amount is too small to warrant concern about its stealing your air.
\end{sikaoti}
